% ===== PRÉAMBULE CANONIQUE — à coller VERBATIM en tête de CHAQUE .tex =====
\documentclass[11pt,a4paper]{article}
\usepackage[utf8]{inputenc}\usepackage[T1]{fontenc}\usepackage{lmodern}
\usepackage[french]{babel}
\usepackage{amsmath,amssymb,mathtools}\usepackage{icomma}
\usepackage{siunitx}\sisetup{locale=FR}  % \qty{}{}, \si{}, \ang{}
\usepackage{xcolor}\usepackage{colortbl}
\usepackage{tikz}\usepackage{pgfplots}\pgfplotsset{compat=1.18}
\usepackage[siunitx,europeanresistors]{circuitikz} % circuits (élec.)
\usepackage[most]{tcolorbox}\usepackage{enumitem}\usepackage{array,booktabs}
\usepackage[a4paper,margin=2.2cm]{geometry}
\usetikzlibrary{arrows.meta,calc,angles,quotes,positioning,patterns,%
  backgrounds,shapes.geometric,decorations.pathmorphing,decorations.markings,intersections}
\definecolor{primary}{RGB}{222,49,99}\definecolor{primarydark}{RGB}{150,22,55}
\definecolor{defcol}{RGB}{0,114,178}\definecolor{methcol}{RGB}{52,92,148}
\definecolor{excol}{RGB}{0,135,90}\definecolor{attncol}{RGB}{200,40,55}
\tcbset{boxbase/.style={enhanced,breakable,boxrule=0.9pt,arc=2.5pt,
  left=3mm,right=3mm,top=2.3mm,bottom=2.3mm,fonttitle=\bfseries,coltitle=white}}
% --- boîtes TUTORIEL (noms EXACTS, ne pas renommer) ---
\newtcolorbox{cle}[1][]{boxbase,colback=defcol!6,colframe=defcol,
  title={\textsc{À retenir}\ifx\relax#1\relax\relax\else\textmd{ : \itshape #1}\fi}}
\newtcolorbox{methode}[1][]{boxbase,colback=methcol!7,colframe=methcol,sharp corners,
  title={\textsc{Méthode}\ifx\relax#1\relax\relax\else\textmd{ --- \itshape #1}\fi}}
\newtcolorbox{exemplebox}[1][]{boxbase,colback=excol!5,colframe=excol,
  title={\textsc{Exemple}\ifx\relax#1\relax\relax\else\textmd{ : \itshape #1}\fi}}
\newtcolorbox{piege}[1][]{boxbase,colback=attncol!6,colframe=attncol,
  title={\textsc{Piège}\ifx\relax#1\relax\relax\else\textmd{ : \itshape #1}\fi}}
\newtcolorbox{remarque}[1][]{boxbase,colback=black!4,colframe=black!45,
  title={\textsc{Remarque}\ifx\relax#1\relax\relax\else\textmd{ : \itshape #1}\fi}}
% --- boîtes EXERCICE / CORRIGÉ (noms EXACTS, auto-comptées) ---
\newtcolorbox[auto counter]{exo}[1][]{boxbase,colback=primary!4,colframe=primary,
  title={\textsc{Exercice}~\thetcbcounter\ifx\relax#1\relax\relax\else\textmd{ --- \itshape #1}\fi}}
\newtcolorbox[auto counter]{corrige}[1][]{boxbase,colback=excol!4,colframe=excol,
  title={\textsc{Corrigé}~\thetcbcounter\ifx\relax#1\relax\relax\else\textmd{ --- \itshape #1}\fi}}
\newcommand{\vv}[1]{\vec{#1}}\newcommand{\ds}{\displaystyle}
\newcommand{\R}{\mathbb{R}}\newcommand{\N}{\mathbb{N}}\newcommand{\Z}{\mathbb{Z}}
\begin{document}
\sisetup{per-mode=symbol}

\begin{corrige}[deux mesures de la vergence de l'\oe il]
\begin{enumerate}[label=\alph*)]
  \item Repos : objet à l'infini $\Rightarrow f=\qty{16}{\mm}=\qty{0.016}{\m}$,
        $C=\dfrac{1}{0,016}=62{,}5~\delta$.
  \item $u=\qty{17}{\cm}=\qty{0.17}{\m}$, $v=\qty{17}{\mm}=\qty{0.017}{\m}$ :
        \[ C=\frac{1}{u}+\frac{1}{v}=\frac{1}{0,17}+\frac{1}{0,017}=5{,}88+58{,}8=64{,}7~\delta. \]
\end{enumerate}
\end{corrige}

\begin{corrige}[vrai ou faux sur l'\oe il et ses défauts]
\begin{enumerate}[label=\alph*)]
  \item \textbf{Faux} : l'\oe il myope est trop long, mais on le corrige par une lentille
        \emph{divergente}.
  \item \textbf{Vrai} : hypermétropie $\to$ vision floue de près $\to$ verre convergent.
  \item \textbf{Vrai} : presbytie $\to$ accommodation insuffisante $\to$ verre convergent.
  \item \textbf{Vrai} : l'iris est le diaphragme de l'\oe il.
  \item \textbf{Faux} : image en arrière de la rétine $\Rightarrow$ \oe il \emph{hypermétrope}
        (le myope forme l'image en avant).
\end{enumerate}
Affirmations correctes : \textbf{b)}, \textbf{c)} et \textbf{d)}.
\end{corrige}

\begin{corrige}[deux myopes, deux corrections]
Objet à l'infini, image virtuelle au point le plus lointain vu net : $v=-D$, d'où $f=-D$ et $C=1/f$.
\begin{enumerate}[label=\alph*)]
  \item Antoine : $f=-\qty{0.90}{\m}$, $C=\dfrac{1}{-0,90}=-1{,}11~\delta$.\\
        Bruno : $f=-\qty{0.50}{\m}$, $C=\dfrac{1}{-0,50}=-2~\delta$.
  \item \textbf{Bruno} est le plus myope (son point lointain est plus proche) ; son verre a la
        vergence la plus forte en valeur absolue ($2~\delta > 1{,}11~\delta$).
\end{enumerate}
\end{corrige}

\begin{corrige}[lire une ordonnance d'opticien]
$f=1/C$ (en mètres) ; le signe donne la nature et le défaut.
\begin{enumerate}[label=\alph*)]
  \item verre A : $C<0$ \textbf{divergent}, $f=\dfrac{1}{-4}=-\qty{0.25}{\m}=-\qty{25}{\cm}$ :
        corrige une \textbf{myopie}.
  \item verre B : $C>0$ \textbf{convergent}, $f=\dfrac{1}{2,5}=+\qty{0.40}{\m}=+\qty{40}{\cm}$ :
        corrige une \textbf{hypermétropie} (ou une presbytie).
  \item verre C : $C=0$ $\Rightarrow$ verre \textbf{neutre} (afocal) : \oe il normal (emmétrope), pas
        de correction.
\end{enumerate}
\end{corrige}

\end{document}
